
\Problem{Circuit values}

\Statement

Verilog is a hardware description language to model digital circuits
that is standardized by IEEE. Write a program that, given the
description of a digital circuit in a simplified version of Verilog
and a sequence of truth assignments for the inputs to the circuit,
outputs the result of evaluating the circuit on those assignments.


\FigureC{width=10cm}{circuitet}


\Input

Input starts with a sequence of output cable identifiers, followed by
a sequence of input cable identifiers, then a sequence of AND, OR or
NOT gate specifications, and finally a sequence of cases. Each case is
a truth-value assignment to the input cables of the circuit in the
same order as given in the description of the circuit. Each cable of
the circuit is identified by a unique sequence of lower-case letters
and digits. The truth values are \texttt{T} and \texttt{F} for TRUE
and FALSE, respectively.  \medskip

Follow the format of the examples. You may assume that the input
describes a correct circuit; that is, input and output cables are
disjoint, all cables are reachable from the inputs, and the circuit
does not contain cycles (in particular, no input cable is the output
of a gate).

\Output

Output is the sequence of truth values that the circuit generates at
its output cables with each assignment, in the same order as given in
the description of the circuit, one per line.

\clearpage
\Sample

